\documentclass[11pt,a4paper]{article}
\usepackage[margin=1in]{geometry}
\usepackage{helvet}
\renewcommand{\familydefault}{\sfdefault}
\setlength{\parskip}{4pt}
\setlength{\parindent}{0pt}

\begin{document}

\begin{center}
{\large\textbf{Reflective Essay}}\\[2pt]
Aggregating Ethics: Moral Uncertainty and Response Selection in Language Models\\[2pt]
Ahmed Idris
\end{center}

\section*{Analysis of Strengths and Weaknesses}

The decision I defend most readily is the axiom harness. Rather than assume that a prompted judge follows its framework, I tested the claim behaviourally, decomposing each theory into five operational axioms and each axiom into three response pairs matched for length and register. The harness earned its keep by finding real faults: the deontological judge initially forgave beneficial lies, and the utilitarian prompt encoded risk aversion where classical utilitarianism demands risk neutrality. Neither defect could be seen by reading the prompts; both were plain in behaviour. That experience taught me to treat validation as a measurement problem before treating it as a philosophical one.

The divergence measure is the second strength. Because integer scores produced ties everywhere, any tie-breaking convention would have manufactured agreement or disagreement by fiat, so I treated each rule's selection as a winner set and made every reported divergence a strict lower bound that no convention could reduce. Splitting disagreement into two levels, the chosen text and the recommended action, then exposed the structure I find most interesting in the whole study. In a fifth of the dilemmas the rules chose different texts that pointed the same way, the labels recording unanimity while the reasoning diverged underneath. A choice made for robustness ended up carrying the analysis, and I have since come to regard measure design as the highest-leverage stage of an empirical project.

Holding the candidate pool and its scores fixed had a similar payoff, since every rule read identical scores for identical candidates and any difference in outcome was attributable to the aggregation function and nothing else. The distillation chapter extended the same instinct into evaluation, reading the copies against references on both sides: a floor from a regression on word counts, and a ceiling from the teacher re-scoring its own work. That floor was humbling, because prompting the base model earned less correlation than counting words on all three judges. A correlation of 0.84 looked impressive until the floor reached 0.76, and the copies' weak ranking only became legible against a self-agreement ceiling near 0.9. I no longer trust a metric that arrives without references on either side of it.

The weaknesses are at least as instructive. Chief among them, the scoring instrument was too blunt for the discriminations I asked of it: on a 0 to 10 integer rubric each judge drew on roughly four distinct values per dilemma, and the utilitarian baseline was indifferent among several candidates on 77 percent of dilemmas. Although I anticipated a ceiling effect while designing the judges, I still underestimated its consequences. The set-valued analysis managed the problem honestly, but a wider, revalidated rubric would have prevented it.

Of all the weaknesses, the recommendation labels weigh most heavily. My headline action-level figure rests on labels produced by a second language model that I checked only for internal consistency, and although a hand-labelled gold sample was feasible in an afternoon, under deadline pressure I chose disclosure over measurement. The disclosure is honest and appears next to the claims it qualifies. An examiner is still entitled to prefer the number, and next time I would spend the afternoon.

The remaining weaknesses are quieter. Sixteen samples from one 8B generator under one neutral prompt gave a homogeneous pool, so the rules often ranked paraphrases of a single recommendation. Meanwhile the keyword mapping used to select contested dilemmas matched fewer than half of the annotated values, which leaves the pool expected to be contested rather than known to be. Because the Ubuntu judge was predicted to contrast with the other judges on only one non-Ubuntu axiom, it may reward communal vocabulary generically rather than Ubuntu reasoning specifically. I also rewrote two judge prompts after watching them fail, then defended the rewrites with six confirmation pairs written after the prompts froze. Their 90 percent pass rate is directionally reassuring and statistically thin, and I have presented it as no more than that.

Looking across the two lists, the strengths share a habit of distrust, each existing because I declined to take something on faith, a judge's fidelity, a tie-break, a bare correlation. By contrast, the weaknesses mostly trace back to instruments that were cheaper or coarser than the questions deserved, accepted early and paid for late. The lesson is one of sequencing: instrument quality bounds everything downstream, and it is the hardest thing to retrofit.

\section*{Possibilities for Further Work}

First on my list is the constraint that binds everything else. Because the candidate pool came from one small generator under one deliberately neutral prompt, its homogeneity suppressed every divergence rate in the study. A generation procedure designed to produce genuinely opposed courses of action would settle whether the modest headline figure reflects the rules or the sampling. Until it runs, the central numbers are lower bounds twice over, by construction of the measure and of the pool.

Two measurement upgrades follow naturally. A scalar lean would replace the categorical action labels, so that divergence could be graded rather than binary and a firm recommendation distinguished from a hedge pointing the same way. A human reading study would then test whether the justification differences my labels conceal are noticeable and consequential for readers. In 101 of the 105 middle-row dilemmas the rules chose fully disjoint texts carrying one label, and whether that difference matters to a person actually seeking advice is an empirical question the dissertation never asks.

The conditional result suggests a practical direction of its own. Ranking disagreement between judges is cheap to compute, and it predicted where the rules part company, with divergence rising fourfold in the most contested quartile. A deployed system could therefore use such a measure as a router, reserving the full aggregation machinery for the dilemmas where it demonstrably matters and defaulting elsewhere, which would turn a diagnostic finding into an engineering economy.

The axiom harness could also grow into a certification suite that stands independently of aggregation. Prompted framework judges are becoming a commodity, and my distillation results show that score correlation certifies very little about them. A public bank of matched pairs, frozen before any training decision, with margins graded rather than passed and ceilings measured rather than assumed, would give the practice a floor. My fifteen axioms and 51 pairs are a small prototype of that idea.

Finally, the credence question deserves better than the equal weights I gave it. I fixed the credences to isolate the aggregation function, and the sensitivity analysis showed that tilting them moves chosen texts far more than recommended actions. Estimating credences empirically, whether from populations or by elicitation, would connect this machinery to the social choice literature it borrows from. Extending the framework beyond three traditions, and testing an Ubuntu judge in a Bantu language rather than in English, are the versions of this project I would most like to see attempted.

\section*{Work I Would Have Conducted with More Time}

A fine-tuning extension was scoped, prototyped, and dropped. The plan used the aggregation rules to build preference pairs and to train one policy per rule, so that differences in selection would become differences in generation. When the prototype's effect size failed to justify the pipeline it required, I cut it in favour of the distillation study, a decision I still believe was right even though it removed the most ambitious question the project had. With more weeks I would have redesigned the experiment around the dilemmas where the rules genuinely part company, rather than abandoning it, since the contested quartile is exactly where a policy-level difference would have the best chance of surviving training noise.

The human gold sample belongs here as well. Roughly a hundred blind labels, stratified across the four classes, would have converted my largest validity threat into a measured agreement rate, and stratification matters because refusals and the rare residual class would otherwise be swamped by the two action labels. I planned the sample, listed it in future work, and did not do it, at a cost of about two hours. The cheapest checks are the easiest to postpone, precisely because they remain feasible right up until the deadline makes them impossible.

Given more time I would also have trained the distilled judges on rationales rather than bare scores. The copies learned their teachers' answers without the reasons, and the results show exactly where reasons matter: their separation margins collapsed from 5.6 points to 1.5 on the short axiom texts, and their framework behaviour barely improved on what the base model already knew. Distilling the score without the justification produced students that imitate rather than understand. Rationale supervision is the obvious next attempt.

A wider scoring rubric sits in the same category, its value obvious from the tie rates by the middle of the project. Because a new rubric is a new instrument, the axiom harness would have had to be rerun to revalidate the judges on it, and that is days of work and API budget I could not spend alongside the distillation chapter.

Budget and compute shaped more of the design than I admitted at the outset. The judges are gpt-4o-mini because scoring 8,000 responses three ways is expensive at any larger scale, and the students are 1.5B models because that is what my hardware trains. A vendor cap of 10,000 requests per day then constrained the re-scoring that set the teacher ceiling. None of these choices invalidates the results, but each narrows the claims. With more resources the first thing I would test is whether the copies' ranking gap is a capacity limit or an information limit, by scaling the student before changing the training data.

\section*{Critical Analysis of the Relationship Between Theory and Practice}

On paper the four rules differ by one line of algebra each. My dissertation exists because that line is misleading: wherever the theory was silent, the implementation had to speak, and the most interesting content of the project lives in exactly those places.

The intertheoretic value problem arrived first, disguised as a preprocessing decision. Although the formal framework assumes choiceworthiness expressed in a common unit, the theories supply no such unit, and in my pipeline that abstract impossibility became a concrete question about standard deviations. Normalising within each dilemma would declare a judge decisive among candidates it had scored between 8 and 9, so I normalised globally instead, committing the study to a broad structural view of how theories compare. A page of philosophical literature turned into four lines of z-scoring code, and the code takes a position in that literature whether I wanted one or not. Writing it forced me to understand the normalisation debate better than reading it ever had.

The Nash rule taught the same lesson from the opposite side. Bargaining theory requires a disagreement point and leaves its location to the modeller, so I fixed it by fiat at a raw score of zero, reading zero as the absence of moral worth on the rubric. The algebra then simplified each judge's surplus to its raw score over its own spread, which is elegant and entirely contingent on my fiat, and a different disagreement point would have crowned different winners across the corpus. Rules that look parameter-free on paper hide their parameters in their inputs, and implementation is where the hiding stops.

Maximin produced the result I could not have predicted from the formalism. Worst-case protection sounds cautious and humane, yet in practice the response whose worst score is highest is often one that no judge scores well, and the flattest score profiles in my pool belonged to safety refusals. On the dilemma about the alcoholic relative, maximin crowned a refusal scored 5, 6 and 5 over every piece of substantive advice on offer --- a theoretical property had become an amplifier for a generator's safety training. No amount of staring at the minimum function would have revealed that, because the effect depends on what texts actually populate the option set. A rule's character is realised only relative to a distribution of candidates, and distributions are engineering facts.

The judges themselves blurred the line between theory and artefact. Against its own theory's clear prediction, my utilitarian judge refused to reward welfare-improving deception on the honesty axiom in every run, because the model beneath the prompt carries an honesty guardrail from alignment training and the guardrail leaked through the consequentialist instruction. The theory assumes an evaluator that simply is the theory; practice supplied an evaluator with a history. I now read every claim about steering a model toward a moral stance with that episode in mind, since the prompt steered the surface while the training held the depths. The same leakage appeared upstream, where the generator refused outright on sensitive dilemmas and shrank the option set before any rule could apply. Moral uncertainty theory takes the option set as given, but in a language model pipeline the option set is itself model behaviour, shaped by the same alignment processes the dissertation set out to interrogate.

Threaded through all of this is the claim I would defend in a viva: aggregation rules that appear interchangeable in the formalism come apart only when implementation forces every free parameter to be set. Expected choiceworthiness and the Nash bargain almost never disagreed about actions once both were given the same normalised inputs, and maximin's distinct character surfaced only through the refusals in the pool. Nor did the traffic run one way, since engineering questions kept sending me back to the normalisation literature while theory kept disciplining the engineering, most visibly in the refusal to break ties that the formalism itself does not break. Reading into moral philosophy and social choice theory as an outsider, then translating what I found into working code, was the most intellectually enjoyable part of the project. Even the two-level divergence measure began as philosophy, since ethics distinguishes the verdict from the reasons for it, and that distinction became the operational split between the recommended action and the chosen text.

My working process taught a parallel version of the same lesson. Throughout the project I used generative AI for code and for drafting, and I learnt to vet it intensely. It invented plausible references that survived until checked against the sources, and more than once it offered a confident suggestion that would have derailed weeks of work had I followed it. So I spent the summer building behavioural tests to catch a model producing plausible imitations of ethical judgement, while working beside a model whose plausible output needed exactly the same scrutiny. Fluency and reliability are different properties. The dissertation argues this about judges, and the process of writing it kept demonstrating the point about assistants.

\section*{Legal, Social, Ethical and Sustainability Issues}

Clearest among the legal and scientific exposures is the dependence on a closed commercial model. Because the three judges are prompts over gpt-4o-mini, the study's central instrument sits behind an API whose weights can change or be withdrawn without notice. Reproducibility therefore decays at the vendor's discretion, and nobody outside the company can audit the model the way my axiom bank audits its behaviour. Distillation was my partial answer, and its outcome sharpens the issue rather than settling it: the copies reproduce the aggregate disagreement rates and fail as judges on closer inspection, so open weights bought portability without buying fidelity. Anyone tempted to publish certified open replacements for closed judges should read that result as a warning. On licensing, the project stacked a research corpus, two open model licences and a commercial API's terms of use, so even a study this small owed compliance to four parties.

The social question is whose ethics get encoded. Against a literature dominated by two Western traditions I included Ubuntu deliberately, and I stand by the choice, though it has a limit I want to name plainly. My Ubuntu judge operationalises one philosopher's reconstruction of a living communal tradition, prompted in English into an American company's model, and the axiom results cannot rule out that it rewards communal vocabulary rather than Ubuntu reasoning. Representation without consultation is thin representation. A better version of this work would involve scholars and speakers of the tradition in authoring the axioms, and would test the judge in a Bantu language. Until then, the inclusion widens the frame while inheriting the very anglophone infrastructure it was meant to balance.

The ethical stakes are concrete because the deployment already exists. People ask language models for advice on precisely the dilemmas in my corpus, yet the selection layer I studied is invisible to those users, and nothing in a chat window discloses whether a cautious worst-case rule or a utilitarian default chose the words on screen. In this study the choice of rule changed the recommended action in about one dilemma in nine, and in more than a quarter of the most contested cases, so a provider's aggregation choice is a moral decision made on users' behalf, usually by accident rather than by design. The maximin result adds a sharper worry, since a worst-case rule systematically elevates refusals and a cautious-sounding design choice can quietly replace advice with disclaimers for the people who most need the former. A combining rule will operate either way, so it should be an explicit, tested and stated design choice.

Sustainability enters through the plain arithmetic of the pipeline. Between scoring, labelling and generation, the study consumed tens of thousands of API calls and 8,000 local generations. The distillation chapter carries the counterweight, since a 1.5B judge runs on consumer hardware, cuts the energy and cost of every subsequent score, and amortises its one-off training cost quickly at any realistic scale. Yet the catch is the chapter's own finding, that the cheap judge is also the unfaithful one. Efficiency and fidelity pulled in opposite directions throughout my results; an honest sustainability argument for small local judges has to carry that tension rather than hide it.

\section*{Closing Remarks}

I have thoroughly enjoyed this project, more than any other piece of work in the degree. Finding a novel problem, formalising it, planning a procedure, and predicting bottlenecks months in advance was formative, and being wrong about some predictions was more formative still. Because model distillation sat outside the syllabus and came from papers, the work drew on most of what the year taught me and pushed past it. Above all it gave me a genuine taste of research, with open questions, homemade instruments, and answers allowed to be modest.

\end{document}
